\documentclass[10pt,a4paper]{amsart}
\usepackage[margin=19mm]{geometry}
\usepackage{amsmath,amssymb,mathtools}
\usepackage[scr=rsfs,cal=euler]{mathalfa}
\usepackage{xfrac}
\usepackage[unicode,hidelinks,bookmarks=false]{hyperref}
\newcommand{\ie}{i.e.\ }
\newcommand{\cf}{cf.\ }
\newcommand{\resp}{respectively\ }
\newcommand{\Subsection}{\subsection}
\providecommand{\nobreakdash}{\nobreak\hbox{-}}
\setlength{\emergencystretch}{2em}
\multlinegap0pt
\usepackage{amssymb}
\usepackage{xcolor}
\usepackage{algorithm}\usepackage{algpseudocode}
\newtheorem{proposition}{Proposition}
\newtheorem{theorem}{Theorem}
\title{Further Results on the Quadratic Embedding Constants of Corona Graphs}
\author{Ferdi, Edy Tri Baskoro, Aditya Purwa Santika}
\date{Source: arXiv:2603.14962v1 (CC BY 4.0)}
\begin{document}
\maketitle

\noindent\emph{Source and licence.} Further Results on the Quadratic Embedding Constants of Corona Graphs. Version arXiv:2603.14962v1 (CC BY 4.0), distributed by arXiv under the Creative Commons Attribution 4.0 International licence, http://creativecommons.org/licenses/by/4.0/, as recorded in the arXiv licence metadata for this exact version and shown on the abstract page https://arxiv.org/abs/2603.14962v1. Full licence notice: \texttt{licenses/arxiv-2603.14962-CC-BY-4.0.txt}.\par\smallskip

\noindent\textit{Editorial scope.} Counted unit: the article's theorem theorem (source0.tex lines 1227--1237). The article's complete original proof of it is delivered with it (source0.tex lines 1239--1510, one \texttt{proof} environment of 272 lines). No mathematics is written, repaired or re-derived: the statement and the proof are the article's own lines, in the article's own notation, and no retained line is substituted or re-flowed.  Retained uncounted as the source-local context the proof needs: lines 718--755.  Nothing was omitted from any retained span, so the package carries no omission record.  No retained line contains a citation, so the wrapper carries no bibliography and no bibliography entry.  No retained line contains a figure environment, an image-inclusion command or drawing code, so no image is delivered and the package declares no asset dependency.  Editorial formatting only, in addition: an AMS \texttt{amsart} preamble that restates the article's own declarations and macros used by the retained lines, namely the environments \texttt{proposition}, \texttt{theorem} on separate counters, together with the standard packages those lines need.  The retained environments print in this document as Proposition 1, Theorem 1 in order of appearance, whereas the article prints the same items with the numbering of its own full document.  The complete native source of the article as distributed in the arXiv e-print for this exact version is bundled unchanged as \texttt{sources/196445/source0.tex}.
\begin{proposition}\label{prop:corona graph of G and H}
Let $G=(V_1,E_1)$ be a connected graph with $|V_1|\ge 2$, and let 
$H=(V_2,E_2)$ be a $\kappa$-regular graph on $n=|V_2|$ vertices, where $n\geq 1$ and $\kappa\geq 0$. 
Let $A_H$ be the adjacency matrix of $H$.
Assume that
\begin{equation}\label{eq:condition-spectrum graph of G and H}
-2-\psi_H^{-1}(\mathrm{QEC}(G))
\notin \mathrm{ev}(A_H).
\end{equation}
Then we have
\begin{equation}\label{eq:qec-corona graph of G and H}
\mathrm{QEC}(G\odot H)
=
\max\{
\psi_H^{-1}(\mathrm{QEC}(G)),\gamma_2,\gamma_3
\},
\end{equation}
where
\begin{equation}
\begin{split}\label{eq: psi of G and H}
\psi_H^{-1}(\mathrm{QEC}(G))&=\frac{1}{2}\bigg\{(n+1)\mathrm{QEC}(G)-(\kappa+2)\\
&\qquad+\sqrt{\left((n+1)\mathrm{QEC}(G)-(\kappa+2)\right)^2+4(\kappa+2)\mathrm{QEC}(G)}\bigg\},
\end{split}
\end{equation}
\begin{equation}\label{eq:gamma2-graph of G and H}
\gamma_2=
\begin{cases}
0, & \text{if } \mathrm{QEC}(G)=0 \text{ or } -2\in \mathrm{ev}(A_H),\\
-\infty, & \text{otherwise},
\end{cases}
\end{equation}
and
\begin{equation}\label{eq:gamma3-graph of G and H}
\gamma_3
=
-2-\min \mathrm{ev}(A_H).
\end{equation}
\end{proposition}

\begin{theorem}
Let $G=(V_1,E_1)$ be a connected graph with $|V_1|\ge 2$, and let 
$H=(V_2,E_2)$ be a $\kappa$-regular graph on $n=|V_2|$ vertices, where $n\geq 1$ and $\kappa\geq 0$. Then
\[
\mathrm{QEC}(G\odot H)=\delta_2(D[G\odot H])
\]
if and only if
\[
\mathrm{QEC}(G)=\delta_2(D_G).
\]
\end{theorem}

\begin{proof}
For simplicity, we write $D:=D[G\odot H]$.

\medskip
\noindent
($\Leftarrow$)
Assume that $\mathrm{QEC}(G)=\delta_2(D_G)$.
To prove the assertion, it suffices to determine the eigenvalues of $D$
arising from the spectrum of $H$ and from the interaction between $G$
and $H$. The following claims describe these eigenvalues.

\medskip
\noindent
\textbf{Claim 1.}
Let $H$ be a $\kappa$-regular graph with adjacency eigenvalues
\[
\kappa=\lambda_1 \ge \lambda_2 \ge \cdots \ge \lambda_n .
\]
Then $-2-\lambda_j$ is an eigenvalue of $D[G\odot H]$ for $j=2,\dots,n$.

Since $H$ is $\kappa$-regular, we have
\[
A_H\mathbf{1}=\kappa\mathbf{1},
\]
and the remaining eigenvectors of $A_H$ are orthogonal to $\mathbf{1}$.
Let $x\in\mathbb{R}^n$ be an eigenvector of $A_H$ corresponding to the eigenvalue $\lambda_j$ with $j\ge2$. Then
\[
A_Hx=\lambda_j x,
\qquad
\langle \mathbf{1},x\rangle=0.
\]
Define the vector $f$ by taking $\xi_o=0$ and
$\xi_i=x_i\mathbf{1}$ for all $i\in V_2$. Substituting these into the previously derived expression for $(D-\lambda)f$, we obtain
\begin{align*}
(D-\lambda)f
&=
\sum_{i\in V_2}D_G(x_i\mathbf{1})\otimes
\begin{bmatrix}1\\\mathbf{1}\end{bmatrix}
+
\sum_{i\in V_2}x_i\mathbf{1}\otimes
\begin{bmatrix}0\\-(A_H+2)e_i\end{bmatrix}
-
\sum_{i\in V_2}\lambda x_i\mathbf{1}\otimes
\begin{bmatrix}0\\e_i\end{bmatrix}.
\end{align*}
Since $\sum_{i\in V_2}x_i=0$, the first term vanishes. Hence
\[
(D-\lambda)f
=
\sum_{i\in V_2}x_i\mathbf{1}\otimes
\begin{bmatrix}0\\-(A_H+2+\lambda)e_i\end{bmatrix}.
\]
Thus $(D-\lambda)f=0$ holds whenever
\[
(A_H+2+\lambda)x=0.
\]
Since $A_Hx=\lambda_jx$, this condition is equivalent to
\[
\lambda=-2-\lambda_j.
\]
Therefore $f$ is an eigenvector of $D[G\odot H]$ with eigenvalue $-2-\lambda_j$.

\medskip
\noindent
\textbf{Claim 2.}
Every root of
\begin{align}\label{eq root of lambda=psi}
\lambda^2-((n+1)\delta_2(D_G)-\kappa-2)\lambda-(\kappa+2)\delta_2(D_G)=0
\end{align}
is an eigenvalue of $D[G\odot H]$.

Let $x$ be a unit eigenvector of $D_G$ corresponding to the eigenvalue
$\delta_2(D_G)$ satisfying
\[
D_Gx=\delta_2(D_G)x,
\qquad
\langle \mathbf1,x\rangle=0,
\qquad
\langle x,x\rangle=1,
\]
which exists since $\mathrm{QEC}(G)=\delta_2(D_G)$. Define
\[
\xi_o=\alpha x, \qquad \xi_i=x, \qquad i\in V_2,
\]
where
\[
\alpha=\frac{n\delta_2(D_G)}{\lambda-\delta_2(D_G)}.
\]
Substituting these into the expression for $(D-\lambda)f$, we obtain
\begin{align*}
(D-\lambda)f
&=\alpha D_Gx\otimes 
\begin{bmatrix}1\\\mathbf1\end{bmatrix}
+nD_Gx\otimes 
\begin{bmatrix}1\\\mathbf1\end{bmatrix}
+x\otimes 
\begin{bmatrix}0\\-(A_H+2)\mathbf1\end{bmatrix}
-\lambda\alpha x\otimes 
\begin{bmatrix}1\\0\end{bmatrix}
-\lambda x\otimes 
\begin{bmatrix}0\\\mathbf1\end{bmatrix}.
\end{align*}
Since $D_Gx=\delta_2(D_G)x$ and $H$ is $\kappa$-regular, we have
\[(A_H+2)\mathbf1=(\kappa+2)\mathbf1.\]
Hence,
\begin{align*}
(D-\lambda)f
&=\delta_2(D_G)(\alpha+n)x\otimes
\begin{bmatrix}1\\\mathbf1\end{bmatrix}
-(\kappa+2)x\otimes
\begin{bmatrix}0\\\mathbf1\end{bmatrix}-\lambda\alpha x\otimes
\begin{bmatrix}1\\0\end{bmatrix}
-\lambda x\otimes
\begin{bmatrix}0\\\mathbf1\end{bmatrix}.
\end{align*}
A direct computation shows that $(D-\lambda)f=0$ holds precisely when
$\lambda$ satisfies the quadratic equation \eqref{eq root of lambda=psi}.
Therefore every root of \eqref{eq root of lambda=psi} is an eigenvalue of
$D[G\odot H]$.

By Proposition~\ref{prop:corona graph of G and H}, we have
\[
\mathrm{QEC}(G\odot H)
=
\max\{\psi_H^{-1}(\mathrm{QEC}(G)),\gamma_2,\gamma_3\}.
\]
Since $\mathrm{QEC}(G)=\delta_2(D_G)$, it follows that
\[
\mathrm{QEC}(G\odot H)
=
\max\{\psi_H^{-1}(\delta_2(D_G)),\gamma_2,\gamma_3\}.
\]
By \textbf{Claim~1}, the values $-2-\lambda_j$ $(j=2,\dots,n)$ are eigenvalues of $D[G\odot H]$, and hence $\gamma_2$ and $\gamma_3$ are an eigenvalue of $D[G\odot H]$. 
Moreover, by \textbf{Claim~2}, $\psi_H^{-1}(\delta_2(D_G))$ is also an eigenvalue of $D[G\odot H]$. 
Therefore
\[
\max\{\psi_H^{-1}(\delta_2(D_G)),\gamma_2,\gamma_3\}
\]
is an eigenvalue of $D[G\odot H]$. Since 
\[
\delta_2(D[G\odot H])
\le
\mathrm{QEC}(G\odot H)
<
\delta_1(D[G\odot H]).
\]
Consequently,
\[
\mathrm{QEC}(G\odot H)=\delta_2(D[G\odot H]).
\]

\medskip
\noindent
($\Rightarrow$)
Assume that
\[
\mathrm{QEC}(G\odot H)=\delta_2(D[G\odot H]).
\]
Then there exists a vector
\[f=\xi_o\otimes \begin{bmatrix} 1 \\ 0 \end{bmatrix}
+ \sum_{i\in V_2} \xi_i\otimes \begin{bmatrix} 0 \\ e_i \end{bmatrix},
\qquad
\xi_i \in C(V_1),
\quad i\in \{o\}\cup V_2\,.\] satisfying
\[
(D-\lambda)f=0, \qquad 
\langle \mathbf1,f\rangle=0, \qquad 
\langle f,f\rangle=1,
\]
where $\lambda=\mathrm{QEC}(G\odot H)$.
From the block structure of $D$, this is equivalent to the
system
\begin{align}
&(D_G-J_G-\lambda)\xi_o 
+\sum_{i\in V_2}D_G\xi_i=0,
\label{05eqn:basic equation (11)}
\\
&\lambda\xi_o
=\sum_{j\in V_2} \langle e_i, (A_H+2+\lambda)e_j\rangle \xi_j,
\qquad i\in V_2,
\label{05eqn:basic equation (12)}
\\
&\langle \mathbf{1}, \xi_o\rangle
+\sum_{i\in V_2} \langle \mathbf{1}, \xi_i\rangle=0,
\label{05eqn:basic equation (13)}
\\
&\langle \xi_o, \xi_o\rangle
+\sum_{i\in V_2} \langle \xi_i, \xi_i\rangle=1.
\label{05eqn:basic equation (14)}
\end{align}

First suppose that $A_H+2+\lambda$ is singular.
Then $\lambda=-2-\theta$ for some adjacency eigenvalue
$\theta\in \mathrm{ev}(A_H)$.
Hence $\lambda\in\{\gamma_2,\gamma_3\}$. Next assume that $A_H+2+\lambda$ is nonsingular.
From \eqref{05eqn:basic equation (12)} we obtain
\[
\xi_i
=\lambda
\langle e_i,(A_H+2+\lambda)^{-1}\mathbf1\rangle
\xi_o,
\qquad i\in V_2.
\]
Substituting this into \eqref{05eqn:basic equation (13)} yields
\[
\bigl(
1+\lambda\langle \mathbf1,(A_H+2+\lambda)^{-1}\mathbf1\rangle
\bigr)
J_G\xi_o=0.
\]
If
\[
1+\lambda\langle \mathbf1,(A_H+2+\lambda)^{-1}\mathbf1\rangle=0,
\]
then \eqref{05eqn:basic equation (11)} implies
\[
(J_G+\lambda)\xi_o=0.
\]
Since $\lambda>-1$ and $\lambda\ne 0$, the matrix $J_G+\lambda$ is nonsingular,
and hence $\xi_o=0$.
Consequently $\xi_i=0$ for all $i\in V_2$, which contradicts
\eqref{05eqn:basic equation (14)}.
Therefore
\[
1+\lambda\langle \mathbf1,(A_H+2+\lambda)^{-1}\mathbf1\rangle\neq0,
\]
and hence $J_G\xi_o=0$.
Using \eqref{05eqn:basic equation (11)} and \eqref{05eqn:basic equation (12)},
we obtain
\[
D_G\xi'-\psi_H(\lambda)\xi'=0,
\qquad
\langle \mathbf1,\xi'\rangle=0,
\qquad
\langle \xi',\xi'\rangle=1,
\]
where
\[
\xi'
=
\left(
1+\lambda^2
\|(A_H+2+\lambda)^{-1}\mathbf1\|^2
\right)^{-1/2}
\xi_o
\]
and
\[
\psi_H(\lambda)
=
\frac{\lambda}{
1+\lambda\langle \mathbf1,(A_H+2+\lambda)^{-1}\mathbf1\rangle}.
\]
Thus, we have $\psi_H(\lambda)$ is an eigenvalue of $D_G$ corresponding to a vector orthogonal to
$\mathbf1$.
Hence
\[
\psi_H(\lambda)\le \mathrm{QEC}(G).
\]
Since $\lambda=\mathrm{QEC}(G\odot H)$, we obtain
\[
\psi_H(\lambda)=\mathrm{QEC}(G).
\]
Finally, since
\[
\delta_2(D_G)\le \mathrm{QEC}(G)<\delta_1(D_G),
\]
we conclude that
\[
\mathrm{QEC}(G)=\delta_2(D_G).
\]
\end{proof}

\end{document}
